% Part: sets-functions-relations
% Chapter: functions
% Section: basics

\documentclass[../../../include/open-logic-section]{subfiles}

\begin{document}

\olfileid{sfr}{fun}{bas}
\olsection{Basisbegrippen}

\begin{explain}
Een \emph{functie} beeldt elk !!{element} van een gegeven verzameling af op
één bepaald !!{element} van een (eventueel andere) gegeven verzameling. Zo
bepaalt de bewerking optellen met~$1$ een functie:
elk getal~$n$ wordt afgebeeld op precies één getal~$n+1$.
  
Algemener kunnen functies paren, drietallen enzovoort als argument hebben en
een bepaalde uitvoer opleveren. Veel functies kennen we uit de elementaire
rekenkunde. Optelling en vermenigvuldiging zijn bijvoorbeeld functies: ze
hebben twee getallen als invoer en leveren één getal als uitvoer op.

In deze wiskundige, abstracte betekenis is een functie een \emph{zwarte
doos}: alleen de koppeling tussen invoer en uitvoer is van belang, niet de
methode waarmee de uitvoer wordt berekend.
\end{explain}

\begin{defn}[Functie]
Een \emph{functie} $f \colon A \to B$ beeldt elk !!{element} van~$A$ af
op een !!{element} van~$B$.

We noemen $A$ het \emph{domein} van~$f$ en $B$ het \emph{codomein}
van~$f$. De !!{element}en van~$A$ heten de invoeren of \emph{argumenten}
van~$f$. Het !!{element} van~$B$ dat aan een argument~$x$ door~$f$ wordt
gekoppeld, heet de \emph{waarde van~$f$} voor argument~$x$ en wordt
geschreven als~$f(x)$.

Het \emph{beeld} $\ran{f}$ van~$f$ is de deelverzameling van het codomein
die bestaat uit de waarden van~$f$ voor een of meer argumenten; $\ran{f} =
\Setabs{f(x)}{x \in A}$.
\end{defn}

Het diagram in \olref{fig:function} kan helpen om functies te begrijpen. De
ellips links stelt het \emph{domein} van de functie voor en de ellips rechts
haar \emph{codomein}. Een pijl loopt van een \emph{argument} in het domein
naar de bijbehorende \emph{waarde} in het codomein.

\begin{figure}
  \olasset{assets/diagrams/function.tikz}
  \caption{Een functie beeldt elk !!{element} van de ene verzameling af op
    !!a{element} van een andere. Een pijl loopt van een argument in het
    domein naar de bijbehorende waarde in het codomein.}
  \ollabel{fig:function}
\end{figure}

\begin{ex}
Vermenigvuldiging neemt paren natuurlijke getallen als invoer en beeldt ze
af op natuurlijke getallen als uitvoer. Dit is dus een functie van
$\Nat \times \Nat$ (het domein) naar $\Nat$ (het codomein). Het beeld blijkt
ook $\Nat$ te zijn, want elke $n \in \Nat$ is gelijk aan $n \times 1$.
\end{ex}

\begin{ex}
Vermenigvuldiging is een functie, omdat zij elke invoer---elk paar natuurlijke
getallen---aan precies één uitvoer koppelt: $\times \colon \Nat^2 \to
\Nat$. Daarentegen levert de koppeling van een natuurlijk getal~$n$ aan een
reëel getal~$x$ waarvoor $x^2 = n$ geen functie op, want elk positief geheel
getal~$n$ heeft twee vierkantswortels: $\sqrt{n}$ en~$-\sqrt{n}$. We krijgen
wel een functie als we alleen de niet-negatieve vierkantswortel kiezen:
$\sqrt{\phantom{X}} \colon \Nat \to \Real$.
\end{ex}

\begin{ex}
De relatie die elke leerling in een klas aan diens eindcijfer koppelt, is een
functie: geen leerling kan binnen dezelfde klas twee verschillende eindcijfers
krijgen. De relatie die elke leerling in een klas aan diens ouders koppelt, is
geen functie: een leerling kan nul, twee of meer ouders hebben.
\end{ex}

\begin{explain}
We kunnen een functie definiëren door voor elk mogelijk argument nauwkeurig
vast te leggen wat haar waarde is. Dat kan met een formule, met een methode om
de waarde te berekenen of met een lijst van de waarden voor alle argumenten.
Welke definitiewijze we ook kiezen, voor elk argument moeten we precies één
waarde vastleggen.
\end{explain}


\begin{ex}
Definieer $f \colon \Nat \to \Nat$ door $f(x) = x+1$. Deze definitie legt
$f$ vast als een functie met natuurlijke getallen als invoer en als uitvoer.
Bij een natuurlijk getal~$x$ levert $f$ de opvolger~$x+1$.
In dit geval is het codomein $\Nat$ niet het beeld van~$f$, want het
natuurlijke getal~$0$ is van geen enkel natuurlijk getal de opvolger. Het
beeld van~$f$ is de verzameling van alle positieve gehele getallen,
$\Int^{+}$.
\end{ex}

\begin{ex}\ollabel{examplefunext}
Definieer $g \colon \Nat \to \Nat$ door $g(x) = x+2-1$. Hieruit volgt dat
$g$ een functie is met natuurlijke getallen als invoer en als uitvoer. Bij een
natuurlijk getal~$x$ levert $g$ de voorganger van de opvolger van de opvolger
van~$x$, dus $x+1$.
\end{ex}

\begin{explain}
We hebben zojuist twee functies, $f$ en $g$, met verschillende
\emph{definities} bekeken. Toch zijn ze \emph{dezelfde functie}. Voor elk
natuurlijk getal~$n$ geldt immers $f(n) = n+1 = n+2-1 =
g(n)$. Anders gezegd: onze definities van $f$ en~$g$ leggen met verschillende
vergelijkingen dezelfde afbeelding vast. We gebruiken hier dus stilzwijgend
het extensionaliteitsbeginsel voor functies:
\[
  \text{als }\forall x\, f(x) = g(x)\text{, dan }f = g
\]
mits $f$ en~$g$ hetzelfde domein én hetzelfde codomein hebben.
\end{explain}

\begin{ex}
We kunnen functies ook gevalsgewijs definiëren. Zo kunnen we
$h \colon \Nat \to \Nat$ definiëren door
\[
h(x) =
\begin{cases}
  \frac{x}{2} & \text{als $x$ even is} \\
  \frac{x+1}{2} & \text{als $x$ oneven is.}
\end{cases}
\]
Omdat elk natuurlijk getal even of oneven is, is de uitvoer van deze functie
altijd een natuurlijk getal. Bij een gevalsgewijze definitie moet elke
mogelijke invoer in precies één geval vallen. Soms moet worden bewezen dat de
gevallen uitputtend en onderling uitsluitend zijn.
\end{ex}
\end{document}
